% GENERATED FILE. Edit canonical lessons, then run npm run generate:books.
% canonical-source-hash: fnv1a64:a2131a2bad2da2f9
% canonical-lessons: IT-C497-educare, IT-C497-viziare, IT-C497-viziato, IT-C497-rimproverare, IT-C497-castigo

\chapter{Crescere i figli}
\label{ch:it-crescere-i-figli}
\hlchaptermodality{497}

\begin{chapteropening}
\textbf{By the end of this chapter:} \emph{I can talk about bringing up children: manners, spoiling and telling off.}

Complete the last lesson of chapter 497: I can talk about bringing up children: manners, spoiling and telling off.
\end{chapteropening}

\section[educare]{educare — to bring up, to educate}

\label{lesson:IT-C497-educare}



\begin{quote}
Before the new one: say the Italian for to rock, then the Italian for to crawl.
\end{quote}

\subsection*{You'll want to know: educare}
\textbf{educare} — \textquotedblleft{}to bring up, to educate\textquotedblright{}.

\begin{quote}
I miei genitori mi hanno educato al rispetto.
\end{quote}

\begin{quote}
\emph{My parents brought me up to be respectful.}
\end{quote}

Mostly about manners and values; for school learning, say \emph{istruire}.

\begin{cousinweb}
\emph{educare} is a learned borrowing from Latin \emph{ēducāre} \textquotedblleft{}to bring up, to rear\textquotedblright{} (EtymDB; Lewis \& Short).
\end{cousinweb}

\begin{grammarlens}[title={What you've built}]
One more word for telling, explaining and describing at B1.
\end{grammarlens}

\subsection*{Your turn}
\begin{itemize}
\raggedright
  \item \emph{Say it:} \emph{educare}
  \item \emph{Say it:} \emph{educare}, once more
  \item \emph{Say it:} a sentence of your own with \emph{educare}
  \item \emph{Recall:} say the Italian for to rock, then the Italian for to crawl, then say \emph{educare} again
\end{itemize}

\begin{tcolorbox}[breakable,colback=teal!4,colframe=teal!35!black,title={Before you move on}]
What does educare mean? (\textquotedblleft{}To bring up, to educate\textquotedblright{}.) Say it once more.
\end{tcolorbox}
\section[viziare]{viziare — to spoil (a child)}

\label{lesson:IT-C497-viziare}



\begin{quote}
Before the new one: say the Italian for to crawl, then the Italian for to bring up.
\end{quote}

\subsection*{You'll want to know: viziare}
\textbf{viziare} — \textquotedblleft{}to spoil (a child)\textquotedblright{}.

\begin{quote}
I nonni viziano sempre i nipoti.
\end{quote}

\begin{quote}
\emph{Grandparents always spoil their grandchildren.}
\end{quote}

Built on \emph{il vizio} \textquotedblleft{}vice, bad habit\textquotedblright{}.

\begin{cousinweb}
\emph{viziare} is a learned borrowing from Latin \emph{vitiāre} \textquotedblleft{}to spoil, to damage\textquotedblright{} (EtymDB), from \emph{vitium} \textquotedblleft{}a fault\textquotedblright{}.
\end{cousinweb}

\begin{grammarlens}[title={What you've built}]
One more word for telling, explaining and describing at B1.
\end{grammarlens}

\subsection*{Your turn}
\begin{itemize}
\raggedright
  \item \emph{Say it:} \emph{viziare}
  \item \emph{Say it:} \emph{viziare}, once more
  \item \emph{Say it:} a sentence of your own with \emph{viziare}
  \item \emph{Recall:} say the Italian for to crawl, then the Italian for to bring up, then say \emph{viziare} again
\end{itemize}

\begin{tcolorbox}[breakable,colback=teal!4,colframe=teal!35!black,title={Before you move on}]
What does viziare mean? (\textquotedblleft{}To spoil (a child)\textquotedblright{}.) Say it once more.
\end{tcolorbox}
\section[viziato]{viziato — spoilt}

\label{lesson:IT-C497-viziato}



\begin{quote}
Before the new one: say the Italian for to bring up, then the Italian for to spoil.
\end{quote}

\subsection*{You'll want to know: viziato}
\textbf{viziato} — \textquotedblleft{}spoilt\textquotedblright{}.

\begin{quote}
È un bambino viziato: vuole tutto subito.
\end{quote}

\begin{quote}
\emph{He is a spoilt child: he wants everything at once.}
\end{quote}

The feminine is \emph{viziata}.

\begin{cousinweb}
\emph{viziato} is the past participle of \emph{viziare} \textquotedblleft{}to spoil\textquotedblright{}, from Latin \emph{vitiāre} (EtymDB), built on \emph{vitium} \textquotedblleft{}a fault\textquotedblright{}.
\end{cousinweb}

\begin{grammarlens}[title={What you've built}]
One more word for telling, explaining and describing at B1.
\end{grammarlens}

\subsection*{Your turn}
\begin{itemize}
\raggedright
  \item \emph{Say it:} \emph{viziato}
  \item \emph{Say it:} \emph{viziato}, once more
  \item \emph{Say it:} a sentence of your own with \emph{viziato}
  \item \emph{Recall:} say the Italian for to bring up, then the Italian for to spoil, then say \emph{viziato} again
\end{itemize}

\begin{tcolorbox}[breakable,colback=teal!4,colframe=teal!35!black,title={Before you move on}]
What does viziato mean? (\textquotedblleft{}Spoilt\textquotedblright{}.) Say it once more.
\end{tcolorbox}
\section[rimproverare]{rimproverare — to scold, to tell off}

\label{lesson:IT-C497-rimproverare}



\begin{quote}
Before the new one: say the Italian for to spoil, then the Italian for spoilt.
\end{quote}

\subsection*{You'll want to know: rimproverare}
\textbf{rimproverare} — \textquotedblleft{}to scold, to tell off\textquotedblright{}.

\begin{quote}
La maestra mi ha rimproverato perché parlavo in classe.
\end{quote}

\begin{quote}
\emph{The teacher told me off because I was talking in class.}
\end{quote}

\emph{Rimproverare qualcuno per qualcosa}. The noun is \emph{il rimprovero}.

\begin{cousinweb}
\emph{rimproverare} is usually traced to Vulgar Latin \emph{reimproperāre}, built on Late Latin \emph{improperāre} \textquotedblleft{}to taunt\textquotedblright{} (Lewis \& Short).
\end{cousinweb}

\begin{grammarlens}[title={What you've built}]
One more word for telling, explaining and describing at B1.
\end{grammarlens}

\subsection*{Your turn}
\begin{itemize}
\raggedright
  \item \emph{Say it:} \emph{rimproverare}
  \item \emph{Say it:} \emph{rimproverare}, once more
  \item \emph{Say it:} a sentence of your own with \emph{rimproverare}
  \item \emph{Recall:} say the Italian for to spoil, then the Italian for spoilt, then say \emph{rimproverare} again
\end{itemize}

\begin{tcolorbox}[breakable,colback=teal!4,colframe=teal!35!black,title={Before you move on}]
What does rimproverare mean? (\textquotedblleft{}To scold, to tell off\textquotedblright{}.) Say it once more.
\end{tcolorbox}
\section[il castigo]{il castigo — punishment}

\label{lesson:IT-C497-castigo}



\begin{quote}
Before the new one: say the Italian for spoilt, then the Italian for to scold.
\end{quote}

\subsection*{You'll want to know: il castigo}
\textbf{il castigo} — \textquotedblleft{}punishment\textquotedblright{}.

\begin{quote}
Per castigo stasera niente televisione.
\end{quote}

\begin{quote}
\emph{As a punishment, no television tonight.}
\end{quote}

\emph{Mettere in castigo} — to punish (a child).

\begin{cousinweb}
\emph{castigo} comes from \emph{castigare} \textquotedblleft{}to punish\textquotedblright{}, a learned borrowing from Latin \emph{castīgāre} \textquotedblleft{}to correct, to chastise\textquotedblright{} (Lewis \& Short), from \emph{castus} \textquotedblleft{}pure\textquotedblright{}.
\end{cousinweb}

\begin{grammarlens}[title={What you've built}]
One more word for telling, explaining and describing at B1.
\end{grammarlens}

\subsection*{Your turn}
\begin{itemize}
\raggedright
  \item \emph{Say it:} \emph{il castigo}
  \item \emph{Say it:} \emph{il castigo}, once more
  \item \emph{Say it:} a sentence of your own with \emph{il castigo}
  \item \emph{Recall:} say the Italian for spoilt, then the Italian for to scold, then say \emph{il castigo} again
\end{itemize}

\begin{tcolorbox}[breakable,colback=teal!4,colframe=teal!35!black,title={Before you move on}]
What does il castigo mean? (\textquotedblleft{}Punishment\textquotedblright{}.) Say it once more.
\end{tcolorbox}
